%% The following is a directive for TeXShop to indicate the main file
%%!TEX root = diss.tex

%=======================================
\chapter{Operators}
\label{ch:Operators}
%=======================================

In this chapter we'll explore some of the properties of the operators that interest us, showing when the commute, if they have inverses, what their eigenvalues and eigenvectors are and so on.


%=======================================
\section{Homogeneous Operators}
%=======================================

%=======================================
\subsection{Single 1D Operators}
%=======================================

We start off with the observation that an extraction operator $E$ is just a marking operator $L$ with mark $a = 0$. In other words, we can combine the discussion for $E$ and $L$ by combining the operators into one definition. 
\begin{definition}
For integer $n$ and a number $a$ let
\begin{align}
	\Gamma_n(a) = 1 + (a-1)x^n[x^n] = 1  + bx^n[x^n]\,.
\end{align}
\end{definition}
\begin{lemma} The operators $\Gamma_m(a_1)$ and $\Gamma_n(a_2)$ commute.
\end{lemma}
\begin{proof} Whether the operators mark the same monomial or not, since they act by multiplication by $a_i$ it doesn't matter in which order we apply them. We can also see this computationally. If $m \neq n$
\begin{align}
	\Gamma_n(a_1)\Gamma_m(a_2) &= 1  +b_1x^n[x^n]+b_2x^m[x^m] + b_1b_2x^n[x^n]x^m[x_m] \\
						&= 1  +b_1x^n[x^n]+b_2x^m[x^m]\\
						&= \Gamma_m(a_2)\Gamma_n(a_1)\,.
\end{align}
The case for $m = n$ is analogous.
\end{proof} 
Because of the commutativity we can write a product of such operators in any order we like. We'll usually write them in alphabetical order. Next, we'll find the eigenvalues and eigenfunctions for single operators.

We note further properties of $\Gamma_n(a)$. The first has to do with its eigenvalues and eigenfunctions. If $a = 1$ the operator acts as the identity so every function is an eigenfunction with eigenvalue 1. When $a \neq 1$ we see that the operator keeps all monomials the same except for the one it marks $x^n$ which will be an eigenfunction with eigenvalue $a$. It's worth seeing this proved in a linear algebra style both to illustrate the method and to act as a consistency check.
\begin{prop} The operator $\Gamma_n(a)$ has eigenvalues $a$ and 1 with corresponding eigenfunctions $x^n$ for $a$ and $x^k$, $k \neq 1$ for 1.
\end{prop}
\begin{proof}
Suppose that for some function $u \neq 0$ and $a \neq 1$
\begin{align}
	\Gamma_n(a) u  = \lambda u\,.
\end{align}
We can write $u = u_n x^n + v$ where $v$ has no $x^n$ term. Therefore we find
\begin{align}
	0 = \Gamma_n(a) u - \lambda u = au_nx^n + v - \lambda u_n x^n -\lambda v = (a-\lambda)u_n x^n + (1-\lambda)v\,.
\end{align}
We have two cases.
\begin{itemize}
\item If $v = 0$ then $u_n \neq 0$ and $\lambda = a$. The eigenfunction is $x^n$.
\item If $v \neq 0$ then $\lambda = 1$ and $u_n = 0$. Here every $x^k$ for $k\neq n$ is an eigenfunction.
\end{itemize}
\end{proof}
The second interesting property has to do with taking powers of $\Gamma_n(a)$, including inverses.
\begin{prop} We have
\begin{align}
	\Gamma_n(a)^m =  \Gamma_n(a^m)
\end{align}
which holds true for all integers $m$ when $a \neq 0$ and all non-negative integers when $a =0$. In particular $\Gamma_n(0)^m = \Gamma_n(0)$.
\end{prop}

When $a = 0$ the operator has a non-trivial kernel generated by $x^n$. We can think of minimal polynomials too. For example $\Gamma_n(0)$ is non-trivial so its minimal polynomial is not $X$ or $X-1$. However since it's idempotent, its minimal polynomial must divide $X(X-1)$ and since it's neither of the factors it must be the minimal polynomial. On the other hand $\Gamma_n(a)$ for $a\neq 0$ does not have a minimal polynomial because the powers of $a$ are independent from each other.

%=======================================
\subsection{Compounded 1D Operators}
%=======================================

Naturally we are led to consider compounded operators and their properties. Let $\Gamma = \Gamma_{n_1}(a_1)\ldots\Gamma_{n_m}(a_m)$ where the powers $n_i$ are distinct; otherwise treat any repetitions using commutativity and the identity $\Gamma_n(a)^m = \Gamma_n(a^m)$. Using induction with the identity above gives
\begin{align}
	\Gamma = 1 + b_1x^{n_1}[x^{n_1}]+\cdots + b_mx^{n_m}[x^m]\,.
\end{align}
Letting $u = u_{n_1}x^{n_1}+\cdots+ u_{n_m}x^{n_m} + v$ in the same way as the single operator case and going through the calculation shows that the eigenvalues of $\Gamma$ are
\begin{itemize}
\item $\lambda = a_i$ for $1\leq i \leq m$ with eigenfunction $x^{n_i}$,
\item $\lambda = 1$ with eigenfunctions $x^m$ where $m \neq n_i$ for any $i$.
\end{itemize}

When one or more of the $a_i$'s is zero we have a non-trivial kernel
\begin{align}
	\ker(\Gamma) = \langle x^m | m = n_i, a_i = 0\rangle\,.
\end{align}

%=======================================
\subsection{Multidimensional Operators}
%=======================================
Multidimensional operators are ones composed of one-dimensional operators (whether single or compounded) and they can also be considered compounded. For this section we'll use superscript notation
\begin{align}
	\Gamma_n^{(k)}(a) = 1 + bx_k^n[x_k^n]\,.
\end{align}
Operators working on different dimensions commute with each other
\begin{align}
	\Gamma_n^{(i)}(a_1)\Gamma_m^{(j)}(a_2)&= 1  +b_1x_i^n[x_i^n]+a_2x_j^m[x_j^m] + a_1a_2x_i^nx_j^m[x_i^n x_j^m]\\
	&= \Gamma_m^{(j)}(a_2)\Gamma_n^{(i)}(a_1)\,.
\end{align}
We can write $\Gamma^{(k)} = \Gamma_{n_1}(a_{k,n_1})\ldots\Gamma_{n_r}(a_{k,n_r})$ as before. This allows us to write a multidimensional operator in terms of its constituent pieces in whichever order we prefer. We are of course interested in finding eigenvalues and eigenfunctions for such operators. We could either find those with the same method as before or we could find it via induction.
\begin{prop} All monomials $x_{[d]}^P$ are eigenfunctions of $\Gamma = \Gamma^{(1)}\ldots\Gamma^{(d)}$. Their eigenvalues are determined by the action of $\Gamma^{(k)}$ on $x_k^{P_k}$. That is, the eigenvalues are products of the eigenvalues produced by the single-dimensional operators.
\end{prop}
This helps us characterize the kernel of $\Gamma$.
\begin{corollary} The kernel of $\Gamma$ is generated by those monomials $x_{[d]}^P$ for which $\Gamma^{(k)}x_k^{P_k} = 0$ for some $k$.
\end{corollary}

%=======================================
\section{Inhomogeneous Operators}
%=======================================
 We'll call an operator inhomogeneous if it consists of a homogeneous operator $\Gamma$ and a laurent polynomial (possibly infinite) $S$ and and acts in one of the following ways
 \begin{align}
	\Gamma S(u) = \Gamma(Su)\quad\text{or}\quad S\Gamma u = S\cdot \Gamma u\,.
 \end{align}
This kind of operator is important because it features quite a bit in the funcitonal equations of random walk models. We seek to find the eigenpairs for  inhomogeneous operators. Our first task is to show that these operators are equivalent.
\begin{prop} The two operators have the same eigenvalues with eigenfunctions differing by a multiple of $S$. That is $(\lambda,u)$ is an eigenpair for $\Gamma S$ if and only of $(\lambda, Su)$ is an eigenpair for $S \Gamma$. In particular, if $\lambda = 0$ we have
\begin{align}
	\ker(\Gamma S) = \ol{S} \ker(\Gamma)\,.
\end{align}
\end{prop}
\begin{proof} Note that
\begin{align}
	\Gamma Su = \lambda u \quad \Longleftrightarrow \quad S\Gamma (Su) = \lambda (Su)\,.
\end{align}
Therefore the operators have the same eigenvalues and they have the eigenfunctions $u$ and $Su$ respectively. Now if we pick $\lambda = 0$ to find the kernels of the operators we find that $u \in \ker(\Gamma S)$ if and only if $Su \in \ker(S\Gamma) = \ker(\Gamma)$ which completes the proof.
\end{proof}

%=======================================
\subsection{Single 1D Operators}
%=======================================
Let $G$ and $u$ be given as
\begin{align}
	S = \sum_k S_k x^k,\quad u = \sum_k u_k x^k
\end{align}
where $u$ is an eigenfunction of $\Gamma S$. We have
\begin{align}
	0&=(\Gamma S-\lambda I) u \\
	&= \sum_{k\neq n} x^k \left(\sum_i S_i u_{k-i} - \lambda u_i\right) + x^n \left(a\sum_i S_i u_{n-i} -\lambda u_n\right)
\end{align}
which leads to the system of equations
\begin{align}
	\lambda u_n &= a\sum_i S_i u_{n-i}\,,\\
	\lambda u_k &= \sum_i S_i u_{k-i}, \quad \text{if } k \neq n\,.
\end{align}
The analysis is non-trivial even for a simple function $S$ like $S = x$. We have
\begin{align}
	\lambda u_n &= a u_{n-1}\,,\\
	\lambda u_k &= u_{k-1}, \quad \text{if } k \neq n\,.
\end{align}
We have a doubly-infinite recursion so we can pick an arbitrary initial condition. If $u_n = 0$ the rest of the terms will be zero as well so we must have $u_n \neq 0$. Set $u_n = 1$. Then via induction we can find
\begin{align}
	u_{n+k} &= 1/\lambda^{k}, \quad \text{for $k \ge 0$}\,,\\
	u_{n-k} &= \lambda^{k}/a^k, \quad \text{for $k \ge 1$}\,.
\end{align}
Computing $u(x)$ formally leads to
\begin{align}
	u(x) &=  \sum_{k=1}^{\infty} u_{n-k}x^{n-k}+\sum_{k=0}^{\infty} u_{n+k}x^{n+k}\\
	        &= x^n\sum_{k=1}^{\infty}\left(\frac{\lambda}{ax}\right)^k +  x^n\sum_{k=0}^{\infty} \left(\frac{x}{\lambda}\right)^k\\
	        &= \frac{x^{n-1}x}{ax-\lambda}-\frac{x^n\lambda}{x-\lambda}\\
	        &=\lambda x^n\left(\frac{1}{\lambda -x}- \frac{1}{\lambda - ax}\right)\\
	        &=-\frac{\lambda b x^{n+1}}{(\lambda-x)(\lambda-ax)}\,.
\end{align}
We could also try using the other definition. Let $u = u_nx^n + v$ where $u$ is now an eigenfunction of $S\Gamma$. Let $S = S_0 + S'$. Then we have
\begin{align}
	S\Gamma u &= (S_0+S')(au_nx^n+v) = \lambda(u_nx^n+v)\\
	0&=(au_nS_0-\lambda u_n+\sum_{k\neq n}S_k u_{n-k})x^n  + au_nS'x^n+S'v-\lambda v \,.
\end{align}
We find that for example
\begin{align}
	(\lambda - a S_0)u_n = \sum_{k\neq n}S_k u_{n-k}
\end{align}
which doesn't seem to help much. 

%=======================================
\section{Random Walks Operator}
%=======================================
In this section we explore the operator $\mc{L} = 1 - t\Gamma S$ for a step Laurent polynomial $S$ and some random walk operator $\Gamma$. We haven't made it rigorous yet but $\mc{L}$ is invertible with inverse
\begin{align}
	{\mc{L}}^{-1} = \sum_{k=0}^{\infty}t^k(\Gamma S)^k
\end{align}
where $(\Gamma S)^k$ is an iteration of the inhomogeneous operator with $(\Gamma S)^0 = 1$. However this inverse is not easy to use as the iteration become quite complicated due to the nature of the random walks. Let $(u, \lambda)$ be an eigenpair for $\mc{L}$.Then we calculate
\begin{align}
	(\mc{L}-\lambda I) u &= u - \lambda u -t\Gamma Su\\
	t \Gamma S u &= (1-\lambda)u = \mu u\,.
\end{align}
If $\lambda = 1$ we seek to find the kernel of $\Gamma S$ which we did in the previous section. The eigenfunctions of $\mc{L}$ for eigenvalue $\lambda = 1$ are the functions $u \in \ol{S}\ker{\Gamma}$.

Next let $u$ be given by a finite order Laurent series
\begin{align}
	u = u_n t^n + u_{n+1}t^{n+1} + \ldots
\end{align}
where $u_n \neq 0$. Then we can re-write the functional equation as
\begin{align}
	t^{-n}(1-t\Gamma S)u = \lambda t^{-n}u\,.
\end{align}
Letting $t \to 0$ in this equation yields
\begin{align}
	(\lambda - 1)u_n = -\lim_{t\to 0} t^{1-n}\Gamma S u = -\Gamma S \lim_{t\to 0} (t^{1-n} u) = 0\,.
\end{align}
Since $u_n \neq 0$ this implies that $\lambda = 1$ and again we see that $u \in \ol{S}\ker \Gamma$. Since $S$ is a Laurent polynomial, the resulting functions will indeed be finite order Laurent series. For the remainder of the analysis we assume that $\lambda \neq 1$. This excludes finite order Laurent series so $u$ is necessarily given by an infinite order Laurent series.


%=======================================
\section{Interaction with $\del_t$}
%=======================================
Here we'll investigate $\del_t$ interacts with $\mc{L}$.
\begin{align}
	\del_t \mc{L} &= \del_t - \del_t(t\Gamma S) = \del_t - \Gamma S- t\Gamma S\del_t = \mc{L}\del_t - \Gamma S\,,\\
	\del_t^2 \mc{L} &= \mc{L}\del_t^2 - 2\del_t \Gamma S\,,\\
	\del_t^n \mc{L} &=\mc{L} \del_t^n - n\del_t^{n-1}\Gamma S
\end{align}
for integer $n$. Another way to write these is with bracket notation
\begin{align}
	[\mc{L},\del_t^n] = \mc{L}\del_t^n-\del_t^n\mc{L} = n\del_t^{n-1}\Gamma S\,.
\end{align}

%=======================================
\section{Interaction with $\del_i$}
%=======================================
Let's take the operator $\Gamma_n(a) = 1 + b x^n[x^n]$ again. Let $\del_k = \del/\del_k$ and in this case just $\del$. We have
\begin{align}
	\del \Gamma_n(a) = \del + b\del(x^n[x^n]) = \del + bnx^{n-1}[x^n]\,.
\end{align}
Now if $n\ge 0$ we have
\begin{align}
	\del^{n+1}\Gamma_n(a)  = \del^{n+1} + b\del^{n+1}(x^n)[x^n] = \del^{n+1}\,.
\end{align}
If $n\leq 0$ taking derivatives will yield ever greater powers of $\ol{x}$ so we want to pre-multiply by a power of $x$ first.
\begin{align}
	\del x^{-n} \Gamma_{n}(a) = \del x^{-n} + b\del [x^n] = \del x^{-n} = -nx^{-n-1} + x^{-n}\del\,.
\end{align}
If we have more terms in an operator then we need to multiply by a power of $x$ to clear out any negative powers of $x$. Then we take the lowest order derivative in order to clear out all the extraction terms. Suppose that
\begin{align}
	\Gamma = 1 + \sum_{i=0}^{m} a_i \ext{x^i} - \sum_{i=1}^n b_i\ext{\ol{x}^i}
\end{align}
where $m$ and $n$ are non-negative. We need to pre-multiply by $x^n$ to clear out all the negative powers of $x$. Then we have
\begin{align}
	x^n\Gamma = x^n + \sum_{i=0}^{m} a_i \ext{x^{i+n}} - \sum_{i=1}^n b_i\ext{x^{n-i}}\,.
\end{align}
This operator is cleared by $\del^{m+n+1}$.
\endinput
